% !Mode:: "TeX:UTF-8"

\section{论文工作是否按计划进行}
请对照开题报告，说明当前总体进度、已完成节点以及是否达到预期要求。若研究内容、技术路线或时间安排发生调整，也请在本节说明原因。

\subsection{总体进度说明}
可从文献调研、理论分析、实验平台搭建、算法实现、系统开发、数据采集与处理等方面进行概述。

\subsection{与开题报告计划的对比}
建议逐项说明原计划、本阶段完成情况、偏差原因以及后续调整方案。

\section{目前已完成的研究工作及阶段性成果}
本节建议围绕“做了什么、怎么做的、取得了什么结果”展开撰写。

\subsection{已完成的研究内容}
可以介绍已经完成的理论推导、模型建立、系统设计、实验方案、程序实现、数据分析等工作。

\subsection{阶段性结果与分析}
可结合图、表、公式或实验结果，说明阶段性成果的有效性，并对结果进行分析和讨论。

\subsection{已有成果沉淀}
如已有论文投稿、专利申请、软件著作权、比赛成果、阶段性技术报告等，可在本节说明。

\section{后续拟完成的研究工作及进度安排}
本节用于说明剩余工作内容、实施路径以及时间规划。

\subsection{后续研究计划}
建议写清楚后续准备完成的关键任务，例如补充实验、完善模型、优化算法、撰写论文各章节等。

\subsection{时间安排}
可按月份或周次给出进度计划，例如：

\begin{center}
  \renewcommand{\arraystretch}{1.3}
  \begin{tabular}{p{0.28\textwidth}p{0.6\textwidth}}
    \hline
    时间阶段 & 计划内容 \\
    \hline
    第1阶段 & 完成补充实验/数据采集，整理阶段性结果。 \\
    第2阶段 & 完成核心方法优化与结果分析，补充对比实验。 \\
    第3阶段 & 完成论文主体撰写与修改，准备送审材料。 \\
    \hline
  \end{tabular}
\end{center}

\section{存在的问题、困难及拟采取的措施}
建议实事求是地说明当前存在的问题，例如实验条件受限、数据不足、方法收敛效果一般、系统稳定性不足等，并给出明确的解决思路。

\section{论文按期完成的可行性分析}
可从前期研究基础、现有实验条件、导师指导、时间安排、经费与数据保障等方面说明按期完成论文的可行性。

\section{参考文献}
\bibliographystyle{hithesis}
\nocite{*}
\bibliography{reference}

% 如学院要求保留“导师意见”“检查小组意见”等栏位，可按需取消注释：
% \section*{导师意见：}
% 请在此处填写导师意见。
%
% \vspace{24bp}
% \hfill 导师签名：\hspace{8\ccwd}
%
% \vspace{24bp}
% \hfill 年\hspace{2\ccwd}月\hspace{2\ccwd}日
%
% \section*{检查小组意见：}
% 请在此处填写检查小组意见。
%
% \vspace{24bp}
% \hfill 组长（签字）：\hspace{8\ccwd}
%
% \vspace{24bp}
% \hfill 年\hspace{2\ccwd}月\hspace{2\ccwd}日

% Local Variables:
% TeX-master: "../master_midterm_harbin"
% TeX-engine: xetex
% End:
